\documentclass[11pt]{amsart}

\usepackage[T1]{fontenc}
\usepackage{lmodern}
\usepackage{microtype}
\usepackage{amsmath,amssymb,amsthm}
\usepackage[hidelinks]{hyperref}

\hypersetup{
  pdftitle={An Order-9 Counterexample to the Triplication-Template Conjecture}
}

\newtheorem{theorem}{Theorem}
\theoremstyle{remark}
\newtheorem*{remark}{Remark}

\newcommand{\Z}{\mathbb Z}
\newcommand{\F}{\mathbb F}
\newcommand{\ST}{\operatorname{ST}}

\title[An order-9 counterexample]
{An Order-9 Counterexample to the Triplication-Template Conjecture}
\date{}

\subjclass[2020]{05B15, 11B75}
\keywords{strong starter, pseudostarter, triplication template,
triplication table, admissible key, modular Sudoku problem, counterexample}

\begin{document}

\begin{abstract}
Ogandzhanyants, Sadov, and Kondratieva conjectured that every
triplication table obtained from a triplication template with an admissible
key is induced by a strong starter. We give a keyed template of order $9$
with an admissible key such that no strong starter in $\Z_{27}$ is congruous
with the resulting table. A short carry identity over $\F_3$ certifies
nonexistence.
\end{abstract}

\maketitle

We use the terminology and numbering of the e-print version of \cite{OSK},
that is, of arXiv:2603.08381v1. For a triplication table $\Sigma_m$, let
$\ST(\Sigma_m)$ denote the set of ordered strong starters in $\Z_{3m}$
congruous with $\Sigma_m$.
Conjecture~7.1 of \cite{OSK} asserts that
\[
 t\in K(T_0,T_1,T_2)
 \quad\Longrightarrow\quad
 \ST\bigl(\Sigma_m(T_0,T_1,T_2,t)\bigr)\ne\varnothing.
\]
Unsolvable triplication tables are already known. Example~7.7 of \cite{OSK}
exhibits two general triplication tables, of orders $11$ and $13$, whose
Modular Sudoku Problem has no solution, and the authors present them as
showing that the extension of Conjecture~7.1 to all triplication tables is
false; Remark~7.8 there adds that such tables are rare under random
sampling, reporting $27$ unsolvable tables among $1000$ generated tables of
order $9$. Neither exhibited table is asserted to arise from a triplication
template. What is new below is thus not an unsolvable triplication table as
such, but one that lies in the template class of \cite[Definition~4.2]{OSK}
with an admissible key, that is, in exactly the class quantified over by
Conjecture~7.1.

\begin{theorem}\label{thm:main}
In $\Z_9$, let
\[
\begin{aligned}
T_0&=[(1,2),(4,6),(5,8),(3,7)],\\
T_1&=[(4,5),(1,3),(8,2),(8,3)],\\
T_2&=[(5,6),(4,6),(7,1),(7,2)].
\end{aligned}
\]
Then $\langle T_0,T_1,T_2\rangle$ is a triplication template,
$1\in K(T_0,T_1,T_2)$, and
\[
\ST\bigl(\Sigma_9(T_0,T_1,T_2,1)\bigr)=\varnothing.
\]
Consequently, \cite[Conjecture~7.1]{OSK} is false.
\end{theorem}

\begin{proof}
The directed differences in each of $T_0,T_1,T_2$, in the displayed
order, are $1,2,3,4$ modulo $9$. The pairing $T_0$ partitions
$\Z_9^*$, while in the multiset union $T_1\uplus T_2$ every element of
$\Z_9^*$ occurs exactly twice. Thus
$\langle T_0,T_1,T_2\rangle$ is a triplication template.

With key $t=1$, the resulting keyed template is
\begin{equation}\label{eq:table}
\Sigma_9=
\begin{array}{c|ccc}
 &0&1&2\\ \hline
0&& (1,1)&\\
1&(1,2)&(5,6)&(6,7)\\
2&(4,6)&(2,4)&(5,7)\\
3&(5,8)&(0,3)&(8,2)\\
4&(3,7)&(0,4)&(8,3).
\end{array}
\end{equation}
Here $0$ occurs twice and every nonzero residue occurs three times. The
three directed differences in regular row $r$ are all $r$, and the multiset
of pair sums modulo $9$ is
\[
\{1,1,1,2,2,2,3,3,3,4,4,4,6\}.
\]
The special pair is $(1,1)$, and the thirteen ordered pairs are distinct
(in fact, their underlying unordered pairs are distinct). Hence
\eqref{eq:table} satisfies Definition~2.2 of \cite{OSK}; equivalently,
$1\in K(T_0,T_1,T_2)$ by Definition~4.3.

Suppose that a strong starter in $\Z_{27}$ is congruous with
\eqref{eq:table}. Index the thirteen table pairs by $i$ and write the
corresponding lifted pair as
\[
 x_i=u_i+9U_i,\qquad y_i=v_i+9V_i,
 \qquad U_i,V_i\in\{0,1,2\}.
\]
Let $r_i\in\{0,1,2,3,4\}$ be the row index, with $r_i=0$ for the special
pair, and let $s_i\in\{0,\ldots,8\}$ be the residue of $u_i+v_i$ modulo
$9$. Define
\[
\delta_i=
\begin{cases}
1,&v_i-u_i<0,\\
0,&v_i-u_i\ge0,
\end{cases}
\qquad
\sigma_i=
\begin{cases}
1,&u_i+v_i\ge9,\\
0,&u_i+v_i<9.
\end{cases}
\]
Then
\begin{align}
 y_i-x_i
 &\equiv r_i+9(V_i-U_i-\delta_i)\pmod{27},
 \label{eq:difference}\\
 x_i+y_i
 &\equiv s_i+9(U_i+V_i+\sigma_i)\pmod{27}.
 \label{eq:addition}
\end{align}
All identities below are in $\F_3$.

Because the lifted pairs partition $\Z_{27}^*$, the three components
above each nonzero residue modulo $9$ have quotient digits $0,1,2$; the
two components above residue $0$ have quotient digits $1,2$. Grouping by
residue therefore gives
\begin{equation}\label{eq:components}
 \sum_i (u_iU_i+v_iV_i)=0.
\end{equation}

For a fixed regular row $r$, put
$D_i=V_i-U_i-\delta_i$. By \eqref{eq:difference}, its three lifted
directed differences are $r+9D_i$. They represent distinct difference
classes. Since $r\not\equiv-r\pmod 9$, two of them cannot be negatives of
one another; hence the three values $D_i$ are $0,1,2$. Summing over the
regular rows yields
\begin{equation}\label{eq:rows}
 \sum_i r_i(V_i-U_i)=\sum_i r_i\delta_i.
\end{equation}

For each $s\in\{1,2,3,4\}$, exactly three pairs in \eqref{eq:table} have
sum $s$ modulo $9$. Their lifted sums are distinct, so by
\eqref{eq:addition} the corresponding values
$U_i+V_i+\sigma_i$ are $0,1,2$. The only other sum class is the singleton
class $s=6$, whose coefficient is zero in $\F_3$. Thus
\begin{equation}\label{eq:sums}
 \sum_i s_i(U_i+V_i)=-\sum_i s_i\sigma_i.
\end{equation}

Add \eqref{eq:components}, \eqref{eq:rows}, and \eqref{eq:sums}. Since
$r_i\equiv v_i-u_i$ and $s_i\equiv u_i+v_i$ modulo $9$, the coefficients
of $U_i$ and $V_i$ are respectively
\[
 u_i-r_i+s_i\equiv3u_i=0,
 \qquad
 v_i+r_i+s_i\equiv3v_i=0.
\]
Every congruous strong starter must therefore satisfy the carry identity
\begin{equation}\label{eq:carry}
 \sum_i r_i\delta_i=\sum_i s_i\sigma_i.
\end{equation}

In \eqref{eq:table}, the only pairs with negative ordinary directed
difference are $(8,2)$ in row $3$ and $(8,3)$ in row $4$. Hence
\[
 \sum_i r_i\delta_i=3+4\equiv1\pmod3.
\]
The summation carries occur for
\[
(5,6),(6,7),(4,6),(5,7),(5,8),(8,2),(3,7),(8,3),
\]
whose sum residues are, respectively,
\[
2,4,1,3,4,1,1,2.
\]
Their total is $18\equiv0\pmod3$, contradicting \eqref{eq:carry}. Thus
$\ST(\Sigma_9)=\varnothing$.
\end{proof}

\begin{remark}
Here $T_0$ is a starter of $\Z_9$, whereas $T_1$ and $T_2$ are
pseudostarters but not starters. Thus the example uses the pseudostarter
generality of \cite[Conjecture~7.1]{OSK}. It leaves open the restriction to
templates built entirely from starters and does not address Horton's
conjecture.
\end{remark}

\begin{thebibliography}{1}

\bibitem{OSK}
O.~Ogandzhanyants, S.~Sadov, and M.~Kondratieva,
\emph{The triplication method for constructing strong starters},
\href{https://arxiv.org/abs/2603.08381v1}{arXiv:2603.08381v1}
[math.CO] (2026); published as
J. Combin. Math. Combin. Comput. \textbf{131} (2026), 115--160,
\href{https://doi.org/10.61091/jcmcc131-08}
{doi:10.61091/jcmcc131-08}.

\end{thebibliography}

\end{document}